\section{The regular pair law in physical heat time}
\label{sec:gt-gap-transport}

Write \(t=\lambda\) for physical heat time throughout the remaining
sections. A dot denotes \(\partial_t\), and a prime on a spatial field
denotes \(\partial_z\). The time variation changes the heat-tilted
measure in \eqref{eq:normalized-field}. The degree remains an independent
parameter.

\subsection{Jointly valid approximation and contour geometry}
\label{sec:gt-physical-phase}

For \(0\le t\le1/2\) the exact error law and the phase estimate take
the form
\begin{equation}
 \begin{aligned}
 \Delta_N(t,z)&:=F_t(z)-G_{N,t}(z)=L_N(t,z)E_{N,t}(z),\\
 L_N(t,z)&=(-1)^N\frac{h_N^*(t)}{(2N)!}z^{2N},
 \qquad |E_{N,t}(z)-1|\le\epsilon_N(|z|,Y).
 \end{aligned}
 \label{eq:gt-physical-error}
\end{equation}
for \(|\Im z|\le Y\le1\), with \(h_N^*=h_N/Z_t=\prod_{j=1}^N\beta_j\)
and the same explicit \(\epsilon_N\) as in \eqref{eq:phase-bound}.
Here the identity holds on the whole complex plane. To see why the
allowance is uniform, perform the Gaussian quadrature and confluent
interpolation of Theorem~\ref{thm:error-law} in \(\mu_t\) at each fixed
time. The monic orthogonality systems are nonsingular since the
density is positive on \(\R\), and all moments and the fixed-degree
interpolation remainder are dominated on compact spatial disks by
the source bound \eqref{eq:tail} with a multiplier
\(\e^{u^2/2+R|u|}\) and a fixed polynomial. In the law of
\eqref{eq:error-random-variable}, the source point has density
\(P_N^2\,\dd\mu_t/h_N^*\); the remaining weights are the same simplex
coordinates as at time zero. The two quantitative inputs in the proof
of Theorem~\ref{thm:phase}, namely
\(\mathbb E_{P_N^2\mu_t/h_N^*}\cosh(4U)\le 2N+12\) and
\(\beta_j+\beta_{j+1}\le B_j\), are already uniform by
\eqref{eq:weighted-cosh} and \eqref{eq:coefficient-envelope}.
Thus \eqref{eq:node-allowances}, \eqref{eq:A-budget} and
\eqref{eq:moment-budgets} apply unchanged. In particular, the real
brackets \eqref{eq:real-sign} hold jointly in \(t\) on their stated
range. The exact variance and fourth moment at degree \(N\) use
links up to \(\beta_{N+1}\) and \(\beta_{N+2}\), respectively, in
\eqref{eq:source-error-moments}. The explicit phase bound uses only
\(N\). At \(t=0\) these are the original formulas.

The signs and value sets in Section~\ref{sec:directed} therefore
apply to each physical snapshot, with its own Jacobi prefix. There
is a distinction between transfer of a boundary count and control
of a ratio logarithm. Write \(\delta(t,z)\) for the separation in the
original field values, namely
\(\operatorname{dist}(-G_{N,t}(z),d_N(z)\mathcal U_N(z))\)
using \eqref{eq:directed-set} at time \(t\).
For \(z\ne0\) it is \(|d_N(z)|\) times the distance in
\eqref{eq:contour-separation}. At zero the permitted error is
\(\{0\}\) and the clearance is \(1\). If that separation is positive
on an open collar of a chosen contour, put \(H_b=G_{N,t}+b\Delta_N\).
The function
\[
 \mathcal L(t,z)=\int_0^1\frac{\Delta_N(t,z)}{H_b(t,z)}\,\dd b
\]
is holomorphic in \(z\) there and its exponential is \(F_t/G_{N,t}\),
by integrating the logarithmic derivative along the parameter \(b\).
The segment of errors \(b\Delta_N\) belongs to the permitted error set,
so \(|\mathcal L|\le|\Delta_N|/\delta\). Uniform clearance on the
chosen collar gives a quantitative logarithm bound. Alternatively,
\(\sup_\Gamma|\Delta_N|<\inf_\Gamma|G_{N,t}|\) gives the bound from the
series for \(\log(1+\Delta_N/G_{N,t})\), which extends to a collar by
strictness. The boundary count and the inner root buffer must each
be justified.

\subsection{Regular drift and transfer of both pair coordinates}

Use the circle \(\Gamma=\{|z-a|=\rho\}\) and pair coordinates from
\eqref{eq:cluster-coordinates}. For each time in a real compact
interval \([u,v]\), suppose \(F_t,G_{N,t}\) are holomorphic on the
closed disk and a neighborhood, obey conjugation symmetry, and
have exactly two counted zeros each with no boundary zeros. For the
transfer estimates below both pairs lie in \(|z-a|\le\rho-b\) with
a common \(b>0\). Suppress time in spatial formulas. Write
\[
 \begin{gathered}
 Q=(z-c)^2-D,\qquad Q_N=(z-c_N)^2-D_N,\\
 A=F_t/Q,\qquad A_N=G_{N,t}/Q_N,\qquad
 \langle f\rangle_\Gamma=\frac{1}{2\pi\ii}\oint_\Gamma f(z)\,\dd z .
 \end{gathered}
\]
Both units are holomorphic and nonzero in the closed disk. Define
\begin{equation}
 K=\left\langle\frac{A'/A}{Q}\right\rangle_\Gamma,\qquad
 K_N=\left\langle\frac{A_N'/A_N}{Q_N}\right\rangle_\Gamma .
 \label{eq:gt-regular-drift}
\end{equation}
These are divided differences of the unit logarithmic derivatives at
the two roots, with the first derivative at a double. If either pair
polynomial is used as the denominator, either unit logarithmic
derivative can be replaced by the corresponding whole-field
logarithmic derivative. The extra rational term has all poles
inside \(\Gamma\) and decays as \(O(z^{-3})\) at infinity.
Conjugation makes the resulting drift real. The formula uses no
division by an internal gap.

\begin{proposition}[Complete pair and drift transfer]
\label{prop:gt-transfer}
Suppose \(|\log(F_t/G_{N,t})|\le\ell_N\) on \(\Gamma\) for a
single-valued holomorphic collar logarithm, and put
\(M_N\ge\sup_\Gamma|G_{N,t}'/G_{N,t}|\). One may take the pointwise
allowances
\begin{align}
 |c-c_N|&\le\eta_{c,N}=\rho\ell_N/2,\qquad
 |D-D_N|\le\eta_{D,N}=\rho(2\rho-b)\ell_N,\notag\\
 |K-K_N|&\le\kappa_N=
 \frac{2\rho\ell_N}{b^3}
 +\frac{\rho M_N}{b^4}\big((4\rho-2b)\eta_{c,N}+\eta_{D,N}\big).
 \label{eq:gt-transfer}
\end{align}
Independently justified smaller coordinate allowances can replace
\(\eta_{c,N},\eta_{D,N}\) in the second line.
The same contour gives \(|K_N|\le\rho M_N/b^2\).
\end{proposition}

\begin{proof}
The proof of Proposition~\ref{prop:pair-resolution} gives the moment
differences \(k\rho^k\ell_N\) by integration by parts. In the
discriminant subtraction, the difference of the center squares is
at most \(2(\rho-b)\eta_{c,N}\) because both centers lie in the
buffered disk. This proves the first line.

Insert the mixed term
\(\langle (G_{N,t}'/G_{N,t})/Q\rangle_\Gamma\).
The first difference is
\[
 \left\langle \frac{(\log(F_t/G_{N,t}))'}{Q}\right\rangle_\Gamma
 =-\left\langle \log(F_t/G_{N,t})(1/Q)'\right\rangle_\Gamma .
\]
The factorization of \(Q\) into two linear terms gives
\(|(1/Q)'|\le2/b^3\). For the second difference, use
\[
 \frac1Q-\frac1{Q_N}=\frac{Q_N-Q}{QQ_N},\qquad
 Q_N-Q=(c-c_N)(2z-c-c_N)+(D-D_N).
\]
On \(\Gamma\), both quadratics have modulus at least \(b^2\), and
\(|2z-c-c_N|\le4\rho-2b\). The circle length is \(2\pi\rho\);
these two bounds give the second line of \eqref{eq:gt-transfer}.
\end{proof}

\begin{proposition}[Physical and perturbed pair laws]
\label{prop:gt-pair-law}
If a field \(Y\) with a counted pair satisfies
\((\partial_t+\partial_z^2+\beta_1)Y=R\), its pair satisfies
\[
 \dot D_Y=2+4K_YD_Y
 -\left\langle (z-c_Y)\frac{R}{Y}\right\rangle_\Gamma .
\]
In particular, the physical field satisfies
\begin{equation}
 \dot D=2+4KD .
 \label{eq:gt-gap-law}
\end{equation}
\end{proposition}
\begin{proof}
Differentiate the two contour moments and integrate the time
derivative of the logarithmic derivative by parts in \(z\).
Their combination in \eqref{eq:cluster-coordinates} gives
\(\dot D_Y=-\langle(z-c_Y)\dot Y/Y\rangle_\Gamma\).
The scalar \(\beta_1\) integrates to zero. For \(Y=Q_YA_Y\), its
second spatial derivative divided by \(Y\) is
\(Q_Y''/Q_Y+2(Q_Y'/Q_Y)(A_Y'/A_Y)+A_Y''/A_Y\).
The last term is holomorphic in the disk. The first contributes
\(\langle Q_Y'/Q_Y\rangle_\Gamma=2\). The cross term uses
\(2(z-c_Y)Q_Y'/Q_Y=4+4D_Y/Q_Y\), whose constant part also
integrates to zero against \(A_Y'/A_Y\). This gives the law
through a double, with \(R=0\) for \eqref{eq:normalized-heat}.
\end{proof}

\subsection{Exact finite Jacobi heat defect}
\label{sec:gt-finite-defect}

Set \(R_N=(\partial_t+\partial_z^2+\beta_1)G_{N,t}\).
For \(N\ge2\), truncate the commutator of \eqref{eq:lax} to the
\(N\)-by-\(N\) matrix, using the strict triangular parts of the
finite \(J_N^2\). The surviving entry on edge \(j\) is
\(\sqrt{\beta_j}(\beta_{j+1}-\beta_{j-1})/2\), with
\(\beta_N=0\) in this finite commutator. Indeed the third
off-diagonal terms involve the same three links and cancel.
The physical \eqref{eq:volterra} restores the last-edge term
\(\beta_N\beta_{N-1}\partial_{\beta_{N-1}}J_N\).
Writing \(f=\e^{\ii zJ_N}\), differentiating in the commutator
direction gives \([B,f]\), where \(B\) is the finite skew part of
\eqref{eq:lax}. Its \(00\)-entry is
\(\sqrt{\beta_1\beta_2}\,f_{20}\) for \(N\ge3\) and zero for \(N=2\).
This cancels the off-diagonal contribution to
\(- (J_N^2f)_{00}+\beta_1 f_{00}\). Hence
\begin{equation}
 R_N=\beta_N\beta_{N-1}\partial_{\beta_{N-1}}G_{N,t}
 =\frac{h_N^*(t)}{2\pi\ii}\oint_\Sigma
      \frac{w\e^{\ii zw}}{P_N(w)^2}\,\dd w .
 \label{eq:gt-residual-contour}
\end{equation}
For the second equality use the spectral resolvent
\(\mathcal R(w)=(w-J_N)^{-1}\) on a contour enclosing its spectrum,
with \(P_N(w)=\det(w-J_N)\). The adjugate entries
\[
 \mathcal R_{0,N-1}
   =\frac{\prod_{j=1}^{N-1}\sqrt{\beta_j}}{P_N(w)},\qquad
 \mathcal R_{0,N-2}
   =\frac{w\prod_{j=1}^{N-2}\sqrt{\beta_j}}{P_N(w)}
\]
come from the unique chain of off-diagonals and the trailing
determinant of size zero or one. Differentiating the resolvent in
the last edge gives
\(\partial_{\beta_{N-1}}\mathcal R_{00}
=\mathcal R_{0,N-2}\mathcal R_{0,N-1}/\sqrt{\beta_{N-1}}\).
This proves the spectral integral. At \(N=1\) only the integral
identity is used and gives \(R_1=\beta_1\), since \(P_1(w)=w\)
and \(G_1=1\).

Let \(V_N\) be the doubled-node simplex barycenter of
\eqref{eq:beta-split}. The confluent exponential formula
\eqref{eq:confluent-exponential}, applied to the \(2N\) doubled
nodes, gives the integral without \(w\) in
\eqref{eq:gt-residual-contour} as
\((\ii z)^{2N-1}\chi_N(z)/(2N-1)!\), where
\(\chi_N(z)=\mathbb E\e^{\ii zV_N}\). Equivalently expand the
integrand at infinity and use \eqref{eq:dirichlet-moments}.
Apply \(\partial_z/\ii\) to insert \(w\). The entire result is
\begin{equation}
 R_N(t,z)=(-1)^{N-1}\frac{h_N^*(t)z^{2N-2}}{(2N-2)!}\Psi_{N,t}(z),
 \qquad \Psi_{N,t}(z)=\chi_N(z)+\frac{z\chi_N'(z)}{2N-1}.
 \label{eq:gt-residual-phase}
\end{equation}
The doubled-node law has fourth moment
\(12(T_2^2+T_4)/(2N)_4\) by the coefficient of degree four in
\eqref{eq:dirichlet-moments}, with \(T_j=\operatorname{tr}J_N^j\).
Subtract its zero mean in the exponential and in its first
derivative. For \(|\Im z|\le Y\le1\), the respective errors are
bounded by \(|z|^2 M_Y/2\) and \(|z|M_Y\) with
\(M_Y=\mathbb E(V_N^2\e^{Y|V_N|})\).
Cauchy--Schwarz, symmetry and \eqref{eq:node-exponential} give
\begin{equation}
 \begin{aligned}
 |\Psi_{N,t}(z)-1|&\le\epsilon_N^R(|z|,Y),\\
 \epsilon_N^R(|z|,Y)&=
 |z|^2\left(\frac12+\frac1{2N-1}\right)
 \left[
 \frac{24((T_N^*)^2+\ell_N^2T_N^*)}{(2N)_4}
 \exp\!\left(\frac{Y^2T_N^*}{N^2(1-Y^2\ell_N^2/N^2)}\right)
 \right]^{1/2}.
 \end{aligned}
 \label{eq:gt-residual-phase-bound}
\end{equation}
Here \(T_N^*,\ell_N\) are as in \eqref{eq:node-allowances}.
The right side is zero at \(N=1\), and is
\(O(|z|^2\log^2(N+1)/N)\) uniformly on the physical slab.
No \(t\)-derivative of the error phase is required.

Proposition~\ref{prop:gt-pair-law} now gives the actual model law
\begin{equation}
 r_N=-\left\langle(z-c_N)\frac{R_N}{G_{N,t}}\right\rangle_\Gamma,
 \qquad \dot D_N=2+4K_ND_N+r_N .
 \label{eq:gt-model-gap}
\end{equation}
The same moment differentiation in the last link gives
\(r_N=\beta_N\beta_{N-1}\partial_{\beta_{N-1}}D_N\) when the
local model count is fixed. At \(N=2\), \(G_2=\cos(\sqrt{\beta_1}z)\)
and \(R_2=-(\beta_2\sqrt{\beta_1}z/2)\sin(\sqrt{\beta_1}z)\).
The residues of \(R_2/G_2\) at its zeros \(z_j\) are \(\beta_2z_j/2\),
so \(r_2=-\beta_2D_2\). This matches the physical variation of the
model roots, which scale as \(\beta_1^{-1/2}\).

\subsection{Signed transport and matched comparisons}
\label{sec:gt-signed-cost}

Put \(W_k(x,y)=\exp(4\int_x^y k(s)\,\dd s)\) for a real integrable
drift \(k\). The physical gap law has the exact forward and backward
forms
\begin{equation}
 D(t)=W_K(u,t)D(u)+2\int_u^t W_K(s,t)\,\dd s,\qquad
 W_K(t,v)D(t)=D(v)-2\int_t^v W_K(s,v)\,\dd s .
 \label{eq:gt-backward-gap}
\end{equation}
In the backward identity an incoming lower value \(L_v\le D(v)\)
strictly exceeding an upper value of the whole cost
\(\mathcal C(u)=2\int_u^v W_K(s,v)\,\dd s\) proves \(D>0\) throughout
\([u,v]\). Equality of the bounds allows \(D(u)=0\) and gives
strict positivity later since the integrand is positive.

Use the model weight \(W_N(t)=W_{K_N}(t,v)\), and define
\[
 \mathcal C_N(x)=2\int_x^v W_N(t)\,\dd t,\qquad
 E_K(t)=4\int_t^v\kappa_N(s)\,\dd s .
\]
The inequalities \(W_N\e^{-E_K}\le W_K(\,\cdot\,,v)\le W_N\e^{E_K}\)
give
\[
 |\mathcal C(x)-\mathcal C_N(x)|
 \le \mathcal T_N^+(x)=2\int_x^v W_N(t)(\e^{E_K(t)}-1)\,\dd t .
\]
Writing \(\mathcal R_N(x)=\int_x^v W_N r_N\,\dd t\) and
\(e_N=D-D_N\), the integrated equations give
\begin{align}
 W_N(x)D_N(x)&=D_N(v)-\mathcal C_N(x)-\mathcal R_N(x),\notag\\
 W_N(x)e_N(x)&=e_N(v)+\mathcal R_N(x)
       -4\int_x^v W_N(t)(K-K_N)(t)D(t)\,\dd t .
 \label{eq:gt-joint-error}
\end{align}
The second follows by subtracting the two differential equations
and differentiating \(W_N e_N\). The last integral has absolute
bound \(4(\rho-b)^2\int_x^v W_N\kappa_N\,\dd t\).
A sufficient strict incoming margin at \(u\) is
\[
 W_N(u)D_N(u)+L_v-D_N(v)+\mathcal R_N(u)-\mathcal T_N^+(u)>0 .
\]
It follows by substituting the first identity into \(L_v-\mathcal C(u)\).
The direct same-data bound remains \(D_N-\eta_{D,N}\).

The weights and signed integrals require continuous-time error
enclosures. For instance, if an evaluator for \(K_N\) has pointwise
error \(\eta_K\) and an evaluated primitive \(\widehat S_N\) has
error \(\eta_S(t)\) relative to the evaluator's exact primitive, then
\[
 s(t)=\eta_S(t)+4\int_t^v\eta_K,\qquad
 \widehat W(t)=\e^{\widehat S_N(t)},\qquad
 |W_N-\widehat W|\le\widehat W(\e^s-1).
\]
A paid upper cost integrand is \(\widehat W\e^{s+E_K}\) and an
evaluator \(\widetilde r_N\) with error \(\eta_r\) satisfies
\[
 |W_Nr_N-\widehat W\widetilde r_N|
 \le \widehat W\big(\e^s\eta_r+(\e^s-1)|\widetilde r_N|\big).
\]
For the contour evaluation of \(r_N\), let the errors in \(c_N,R_N,G_N\)
on the circle be \(\gamma_c,\gamma_R,\gamma_G\). Let \(\widetilde T\)
bound the evaluated numerator, \(m\le\inf_\Gamma|G_N|\) with
\(\gamma_G<m\), and \(g=2\rho-b\). Subtracting the products and
reciprocals in \eqref{eq:gt-model-gap} gives
\[
 \eta_r=
 \rho\left(\frac{g\gamma_R}{m}
 +\frac{g\widetilde T\gamma_G}{m(m-\gamma_G)}
 +\frac{\gamma_c\widetilde T}{m-\gamma_G}\right)+\eta_{\Gamma,r},
\]
where \(\eta_{\Gamma,r}\) encloses the contour integration error.
Each nonnegative error integral needs an upper enclosure and each
evaluated signed integral needs its time quadrature allowance.
Direct endpoint comparison also needs an enclosed zero count and
the endpoint logarithm. For a direct approximation of the normalized
source \(F_t=J_t/Z_t\), let \(J_{\rm app}/Z_{\rm app}\) be the
approximant, with holomorphic numerator and \(Z_{\rm app}>0\).
If the numerator error is \(\epsilon_J\), the mass error is
\(\epsilon_Z\), and \(Z_t\ge Z_{\min}>0\), subtraction gives a
field error bound
\((\epsilon_J+\|J_{\rm app}/Z_{\rm app}\|\epsilon_Z)/Z_{\min}\)
on the contour. The same logarithm comparison then applies.

The residual and \(G_{N+1,t}\) use the same first \(N\) actual
links. On fixed compact subsets of \(|\Im z|\le1\) avoiding
\(z=0\), \eqref{eq:gt-physical-error} and
\eqref{eq:gt-residual-phase} give
\[
 \frac{|R_N(t,z)|}{|F_t(z)-G_{N+1,t}(z)|}
 =\frac{\Pi_N}{\beta_{N+1}(t)|z|^4}
     \frac{|\Psi_{N,t}(z)|}{|E_{N+1,t}(z)|},
 \qquad \Pi_N=(2N-1)(2N)(2N+1)(2N+2).
\]
The last factor tends to one uniformly on these compacts and is
enclosed by the two phase bounds as soon as both are smaller than
one. Thus the pointwise residual is at least of order
\(N^4/\log^2(N+1)\) times that field error. A contour comparison
similarly distinguishes the particular absolute integration allowance
from an actual signed error. Put \(Z_*=\max_\Gamma|z|\),
\(H=\sup_I h_N^*\), \(\bar\beta\ge\sup_I\beta_{N+1}\), and
\[
 T_R=\frac{HZ_*^{2N-2}}{(2N-2)!}(1+\epsilon_N^R(Z_*,1)),\qquad
 \delta_+=\frac{H\bar\beta Z_*^{2N+2}}{(2N+2)!}
               (1+\epsilon_{N+1}(Z_*,1)).
\]
For positive model floors \(m_N,m_{N+1}\) on the same buffered
circle, \(\eta=\delta_+/m_{N+1}<1\) and
\(M\ge\sup_I K_N\), the absolute residual cost is bounded by
\(\rho(2\rho-b)(T_R/m_N)Q_{v-u}(M)\) where
\(Q_h(M)=\int_0^h\e^{4Ms}\,\dd s\). The direct next-head allowance
is \(\rho(2\rho-b)[-\log(1-\eta)]\). Their ratio is
\[
 \frac{m_{N+1}}{m_N}Q_{v-u}(M)
 \frac{\eta}{-\log(1-\eta)}
 \frac{\Pi_N}{\bar\beta Z_*^4}
 \frac{1+\epsilon_N^R(Z_*,1)}{1+\epsilon_{N+1}(Z_*,1)} .
\]
For comparable floors, bounded \(M\), fixed contour and
\(0<\eta\le1/2\), this has the scale
\((v-u)N^4/\log^2(N+1)\) when using \(\bar\beta=B_N\).
It compares one absolute residual allowance with one direct
allowance. The actual signed residual and the full correction
in \eqref{eq:gt-joint-error} have different magnitudes in general.

A drift band \(k_-\le K\le k_+\) is shared information. If
\(0\le L_u\le D(u)\), the forward lower envelope substitutes
\(W_{k_-}\) and \(L_u\) in the forward formula. For the backward
envelope set \(N_+(t)=L_v-2\int_t^v W_{k_+}(s,v)\,\dd s\).
If \(N_+\ge0\), division by the upper weight \(W_{k_+}(t,v)\)
gives a lower bound. If \(N_+<0\), use the positive lower weight
\(W_{k_-}(t,v)\). To see the sign requirement, take the scalar
law with \(D(0)=1\), \(K=0\) on \([-1,0]\) and \(K=-1/4\) on
\([-2,-1)\). Then \(D(-1)=-1\) and \(D(-2)=2-3\e<-3\), while
substituting the upper drift zero in both positions gives \(-3\).
For a constant \(k\) on a cell of oriented length \(h\), the map is
\(y\mapsto \e^{4kh}y+(\e^{4kh}-1)/(2k)\), with additive term
\(2h\) at \(k=0\). For sign-changing comparisons, the forward
lower field uses \(2+4(k_-y_++k_+y_-)\) with \(y_\pm=\max(y,0)\),
\(\min(y,0)\), respectively. In reverse time the lower field is
\(-2-4(k_+y_++k_-y_-)\) with reversed time arguments. These are
the pointwise minima over the drift band and comparison follows
by applying the scalar Lipschitz inequality to the positive
part of the difference. The direct and transport methods can both
compose these cell maps and intersect the bounds with direct
coordinate intervals. For fixed signs this uses a linear number
of scalar map operations in the number of cells. The operation
count leaves the cost of the shared underlying field and geometry
enclosures to their own evaluation.

\begin{proposition}[Minimum on the fixed Lehmer tube]
\label{prop:gt-matched-minimum}
For the fixed tube on \(0\le t\le1/10\) centered at \(a=140101/10\),
the radius bound \(3/4\) and drift band certified in
Appendix~\ref{sec:gt-source-certification},
\(-5/8\le K\le0\) imply \(D'(t)\ge19/32\). The minimum of \(D\)
is the direct endpoint value \(D(0)\).
\end{proposition}
\begin{proof}
The radius bound gives \(|D|\le9/16\).
For \(D\ge0\), \eqref{eq:gt-gap-law} yields
\(D'\ge2-(5/2)(9/16)=19/32\).
For \(D<0\), the product \(KD\) is nonnegative, so \(D'\ge2\).
The already enclosed direct endpoint belongs to
\([0.00142117117778,0.00142118228384]\). The same derivative
inference is available to both methods.
\end{proof}
